\section{Introduction}
\label{sec:intro}

Bi-temporal remote-sensing change detection (CD) asks, given two co-registered
images of the same scene taken at different times, to segment the pixels whose
semantic content has changed. It underpins applications such as urban growth
monitoring, disaster response, deforestation tracking, and infrastructure
auditing~\citep{daudt2018fully,chen2020spatial}. Despite the maturity of the
task, progress in the last two years has been driven almost exclusively by
\emph{specialist} architectures --- Siamese encoders with custom fusion
modules~\citep{fang2021snunet,chen2021remote}, hybrid CNN/transformer
backbones~\citep{bandara2022changeformer,zheng2022changestar}, or diffusion-based
feature extractors~\citep{bandara2024ddpmcd}. Meanwhile, general-purpose vision
has converged on a small family of very strong self-supervised foundation
models: \textsc{DINOv2}~\citep{oquab2023dinov2}, \textsc{DINOv3}~\citep{simeoni2025dinov3},
and the \textsc{SAM} family for segmentation~\citep{kirillov2023segment,ravi2024sam2}.

A natural question follows: \emph{are specialist CD architectures still necessary,
or do frozen foundation models, connected by a lightweight learned adapter,
already contain what a change detector needs?}

Recent work has begun to answer this question positively for individual
foundation-model choices. TTP~\citep{chen2024ttp} adapts \textsc{SAM}'s image
encoder with LoRA and reports strong results on LEVIR-CD.
BAN~\citep{li2024ban} freezes a CLIP backbone and learns a small
change-detection adapter. SAM-CD~\citep{dingsamcd} uses SAM's mask decoder
with a ViT encoder. However, these works (a) fix a single choice of encoder and
decoder family, and (b) typically still require moderate-size learnable
modules ($>$10~M parameters) and careful pre-processing. A systematic study of
foundation-model pairings for CD, across \emph{both} axes (encoder family and
mask-decoder family), is missing.

\paragraph{Contributions.} We present \textbf{GALASH}, a unified
change-detection pipeline and ablation platform. Our contributions are:

\begin{enumerate}
  \item \textbf{A unified encoder/decoder registry} exposing 16 DINO-family
    encoders (DINOv2, DINOv2-with-registers, DINOv2-RS, DINOv3, DINOv3-sat)
    and 8 SAM-family decoders (SAM-1, SAM-2.1 tiny/small/base/large, SAM-3)
    behind a single architecture. To our knowledge this is the first paper to
    systematically compare the 100+ encoder--decoder combinations this
    induces.
  \item \textbf{A Cross-Change Attention module} that uses a learnable
    temperature to produce per-patch change representations from DINO features,
    and which is jointly supervised with (i) the final mask and (ii) a patch-level
    latent change loss.
  \item \textbf{Bridge v2}: a multi-scale fusion block (FPN + transformer
    refinement) that feeds the SAM decoder's \emph{image embedding}, \emph{dense
    prompt}, and \emph{high-resolution} slots. Critically, we repurpose SAM's
    dense-prompt channel --- normally unused in mask-only inference --- as a
    change-conditioning signal.
  \item \textbf{Only $\sim$16.5~M trainable parameters} on top of frozen
    foundation backbones, with competitive results on six standard CD
    benchmarks. This is 2--6$\times$ fewer trainable parameters than recent
    specialist baselines.
\end{enumerate}

\paragraph{Paper outline.} Section~\ref{sec:related} reviews specialist CD
architectures and foundation-model adapters. Section~\ref{sec:method} details
GALASH, the registry, cross-change attention, and Bridge v2.
Section~\ref{sec:experiments} reports per-benchmark comparisons to SOTA;
Section~\ref{sec:ablations} isolates each component. Section~\ref{sec:conclusion}
discusses limitations and directions.
